\section{Final years, martyrdom, and burial}\label{sec:paul-death}

Paul's last well-documented plan and Acts' last scene do not meet in a surviving
execution narrative. Romans looks forward: he has completed a phase of work
from Jerusalem around to Illyricum, intends to carry the nations' collection
to Jerusalem, hopes to visit Rome, and wants Roman support for Spain. Acts
looks back over the collection journey and forward from an unresolved Roman
custody. Early Christian witnesses then remember martyrdom. The gaps between
those horizons are where later lives have most readily become overprecise.

\subsection{Jerusalem and arrest}

Romans asks the recipients to pray that Paul will be rescued from unbelievers
in Judea and that his service will be acceptable to the saints. Acts narrates
his determination to go despite prophetic warnings. At Jerusalem, James and
the elders receive him and propose that he join men completing a vow to answer
reports that he teaches diaspora Jews to abandon Moses. Paul accepts. Jews
from Asia then accuse him of teaching against the people, law, and Temple and
of bringing Greeks into the sacred precinct. Acts says they had seen Trophimus
with him in the city and assumed Paul had brought him inside. A Roman tribune
intervenes as the crowd attacks.

The Acts sequence is theologically dense. Paul addresses the crowd in Hebrew
or Aramaic, recounts his call, claims Roman citizenship when flogging is
prepared, appears before the Sanhedrin, and receives a promise that he will
witness in Rome. A plot leads to transfer under guard to Caesarea. Hearings
before Felix, Festus, and Agrippa display Christianity as connected to Israel's
hope and not a simple civic insurrection. Paul appeals to Caesar; Festus sends
him to Rome.

The letters do not independently narrate these proceedings. Acts' named
officials and local details give the story historical texture, while the
speeches remain Lukan compositions. Paul's appeal depends on the citizenship
Acts attributes to him. It cannot be used to construct the wording of a lost
trial record.

\subsection{Voyage and the open ending of Acts}

Acts 27--28 gives an unusually detailed sea voyage: changing ships, delayed
sailing, storm, shipwreck on Malta, winter, and the journey through Syracuse,
Rhegium, and Puteoli. Roman believers meet Paul at the Forum of Appius and
Three Taverns. Under guarded lodging he explains himself to local Jewish
leaders, insists that he is chained for Israel's hope, and proclaims for two
years to all who come. The final sentence emphasizes unhindered proclamation,
not the prisoner's legal fate.

That ending has been read as evidence that Luke wrote before a verdict, as a
deliberate theological conclusion, or as a sign that Paul was released. The
text alone does not decide among them. It records neither acquittal nor
execution. The traditional language of a ``first Roman imprisonment'' already
assumes a second and should be used only within the hypothesis that Paul was
released and rearrested.

\subsection{Release, Spain, and a second arrest}

Paul intended to go to Spain after Rome (Rom 15:24, 28). Intention is secure;
completion is not. First Clement says that Paul reached the ``limit of the
West'' before bearing final witness. From a Roman author's perspective the
phrase may point to Spain, but it can also function rhetorically or describe
Rome as the western goal of the mission. It is valuable early evidence, not a
passport stamp.

Eusebius reports that after defending himself Paul was sent again to the
ministry of preaching and then returned for martyrdom. The Pastorals imply
travels to Crete, Ephesus,
Macedonia, Miletus, Troas, and Nicopolis not easily fitted before Acts' end and
present Second Timothy as a final Roman testament. Defenders of direct
authorship use them to reconstruct release and a second arrest; critics place
the letters in post-Pauline reception. No combination reaches certainty.

This publication leaves both responsible possibilities open. Paul may have
been released, perhaps traveled west or east, and later returned to Roman
custody. He may instead have died after the custody with which Acts ends. Spain
is a meaningful intended horizon and an ancient received claim, not an
established chapter of the life.

\subsection{The early Roman martyr memory}

First Clement is treated here as the earliest located external witness to
Paul's completed course, although its precise date was not independently
audited in this generation. It remembers
chains, exile, stoning, preaching in east and west, witness before rulers, and
his departure from the world after reaching the boundary of the West. The
chapter pairs him with Peter as an example from the preceding generation. It
does not name Nero, describe a trial, say ``beheaded,'' or give a calendar
date.

Ignatius tells Roman Christians that he does not command them as Peter and
Paul did; this study did not independently audit the letter's precise date or
recensional history. Dionysius of Corinth links both
apostles' teaching in Italy and martyrdom at about the same time; Irenaeus
associates them with the foundation and ordering of the Roman church. Gaius
points to apostolic memorials at the Vatican and on the Ostian Way.
Tertullian explicitly associates Paul's Roman death with beheading. Origen,
as preserved by Eusebius, says Paul suffered martyrdom at Rome under Nero but
does not specify the manner in the checked passage.

This convergence makes Roman martyrdom under Nero probable, commonly in the
range 64--68. Tacitus supplies the hostile context of Nero's punishment of
Christians after the fire of 64 but does not name Paul. The later year 67, a
same-day execution with Peter, a named judge, last words, and a detailed road
to the scaffold are not secured by the earliest witnesses. Beheading is early
received tradition and plausible. No Roman-law authority inspected for this
study links the citizenship attributed by Acts to that manner of death, and no
sentence survives.

\subsection{The Ostian memorial and sarcophagus}

Gaius's report of a Pauline \emph{tropaion} on the Ostian Way is the earliest
located textual pointer to the place later monumentalized by the Basilica of
Saint Paul Outside the Walls. The term can mean memorial or trophy and does not
by itself prove intact remains. A fourth-century basilica was built over the
remembered grave and soon enlarged; despite rebuilding after the 1823 fire,
the sanctuary retained its relation to the ancient memorial. A marble
sarcophagus bearing a Pauline dedicatory inscription lies beneath the papal
altar and is officially venerated as the apostle's tomb.

During modern investigation a small probe entered the sarcophagus. Benedict
XVI announced on 28 June 2009 that minute bone fragments, tested by experts
who did not know their provenance, produced a radiocarbon range consistent
with a person who lived in the first or second century. The probe also found
textile and incense traces. That public announcement is an official account
of a scientific result; this study did not locate a complete peer-reviewed
laboratory report, calibration data, individual profile, or independent chain
of custody sufficient to identify Paul.

The finding is therefore consistent with the ancient tradition and does not
authenticate a named person. The date range is wider than Paul's lifetime;
there is no genetic comparator; the contents' depositional history is not an
unbroken modern forensic record. Official veneration, an ancient fixed cult
site, and personal scientific identification remain three different claims.

\subsection{Execution-site and later relic claims}

Tre Fontane, the Mamertine Prison, and dispersed relic claims were encountered
as research leads, but no exact witness for them was inspected in this
generation. Their local narratives and provenance therefore remain uncollated
and are not published here as established devotional history or as parts of a
reconstructed execution route. The Ostian memorial concerns the received
burial site; it does not authenticate a separate place of custody, execution,
or relic merely by association.

\evidenceaxis{Paul probably died as a martyr at Rome under Nero and was
remembered on the Ostian Way in Gaius's ancient report.}{1 Clem. 5; Ignatius,
\emph{Rom.} 4; Dionysius, Gaius, and Origen in Eusebius; Irenaeus;
Tertullian; the Ostian memorial and basilica.}{Acts ends before a verdict.
Release, Spain, second arrest, exact year, beheading scene, execution site,
particular remains, and same-day death with Peter carry different grades.}
